\documentclass[11pt,a4paper]{article}
\usepackage[margin=2.35cm]{geometry}
\usepackage{amsmath,amssymb,booktabs,longtable,array,xcolor}
\usepackage[colorlinks=true,linkcolor=blue!60!black,urlcolor=blue!60!black]{hyperref}
\newcommand{\Ca}{\textbf{(a)}\ }
\newcommand{\Cc}{\textbf{(c)}\ }
\newcommand{\Cd}{\textbf{(d)}\ }
\newcommand{\dd}{\mathrm d}
\newcommand{\file}[1]{\nolinkurl{#1}}
\newcommand{\kernels}{\file{good_twobody_kernels.py}}
\newcommand{\ct}{\file{good_twobody_counterterms.py}}
\newcommand{\three}{\file{good_twobody_threebody.py}}
\newcommand{\grid}{\file{good_twobody_grid.py}}
\newcommand{\insertcheck}{\file{good_twobody_local_insertion.py}}
\emergencystretch=3em
\title{Good-equation two-soliton scattering at one loop:\\
fixed-momentum counterterm response, three-soliton tests,\\
and the unresolved antisymmetric phase}
\author{Quantum Boussinesq solitons project --- PI: Gustav Delius}
\date{4 November 2026}
\begin{document}
\maketitle
\begin{abstract}
\Cc We independently evaluate the audited continuum integrals in both
channels, verify three-soliton phonons by four-wave Hirota coefficient tests,
and compute the response of the classical canonical scattering phase to
\(\Delta H=\lambda\hbar\int u^2\dd x\) at fixed physical momenta:
\[
 K_\lambda=3\delta_0+4\sum_{j=1}^2P_j\Delta X'_j
 =\frac23(v_2-v_1)\log A_{12}+\text{finite endpoint terms}.
\]
It is antisymmetric under exchange of the complete particle labels and is
generically nonzero. Thus the canonical antisymmetric response is not
scheme independent. We also give responses to finite changes of the kinetic,
cubic and bending coefficients. A separate connected local-insertion
calculation checks the logarithmic term.

\Cc Scalar three-body factorisation and the exact factorised phonon
transmission do not determine an unknown antisymmetric pair contribution.
We exhibit a rank-two pair function that passes these tests; this is a
counterexample to sufficiency of the tested constraints, not a freely
adjustable physical coupling. \Cd The boundary-compatible projected
two- and three-soliton determinants and quantum state map remain
uncomputed. Consequently this is a partial result, not the requested
absolute physical one-loop phase. The JSON grid has 270 rows with explicit
unknown physical entries, numerical ingredients, counterterm coefficients,
and separate near-resonance data. No unknown is replaced by zero.
\end{abstract}

\section{Scope and controlling evidence}
Labels (a), (b), (c), (d) mean literature proof, literature heuristic,
project derivation, and guess or unestablished identification. All new
derived formulas below are checked by the script named in their section;
scripts and caches are in \file{phase2/scripts/}. Exact algebra and
finite numerical tests are distinguished. No legacy script was run or
imported; only \file{good_twobody*} files were written. No renormalisation
condition or quantum boundary prescription is selected.

\Cc The controlling inputs are \file{audit_det_projected.tex},
\file{audit_projected_measure.tex}, \file{audit_determinant_foundations.tex},
\file{audit_smx.tex}, \file{audit_frame_dictionary.tex}, and
\file{errata.tex}; the corresponding package claims are used only with
their audit qualifications. In particular, neither the raw phonon integral
nor its antisymmetrisation has been established as the full projected
determinant. The Toda Krein equation is an additional microscopic input in
a different model. It is not a consequence of scalar factorisation and is
not imported into the good equation. This is a statement of the supplied
audit record, not a new literature theorem.

The fixed convention is
\begin{gather*}
 u_{tt}=u_{xx}-u_{xxxx}-3(u^2)_{xx},\qquad u=\phi_x=2(\log\tau)_{xx},\\
 L=\tfrac12\phi_t^2-\tfrac12u^2+u^3-\tfrac12u_x^2,\qquad
 \Omega=\int\dd\pi\wedge\dd\phi,\quad \pi=\phi_t.
\end{gather*}
All real quantum candidates in the numerical grid are restricted to
\(0<\alpha_j<\pi/3\), \(A_{12}>0\), and unequal velocities.
Angles are in radians. A displayed ordered phase refers to particle 1
overtaking particle 2; exchanged labels denote the inverse oriented
generator. Both \(++\), \(--\), \(+-\), and \(-+\) are saved.

\section{Classical normalization and fixed-momentum conversion}
\paragraph{Checked input \Cc\ (\ct, \kernels).}
With \(k_j=\sin\alpha_j\), \(v_j=\epsilon_j\cos\alpha_j\),
\[
 M_j=2k_j,\quad \Pi_j=-2k_jv_j,\quad P_j=\frac23v_jk_j^3,
 \quad E_j=\frac23k_j^3-\frac8{15}k_j^5,\quad
 \ell_j=-\epsilon_j\cos3\alpha_j.
\]
The energy is the errata-corrected one. Define
\begin{align}
 p_{ja}&=\tfrac12(k_j+v_j/\sqrt3),&
 p_{jb}&=\tfrac12(v_j/\sqrt3-k_j),\\
 A_{12}&=\frac{(p_{1a}-p_{2a})(p_{1b}-p_{2b})}
 {(p_{1a}-p_{2b})(p_{1b}-p_{2a})},&
 T_{12}&=\frac23M_1M_2(v_1-v_2),\\
 \delta_0&=\frac29(\ell_1-\ell_2)\log A_{12}+T_{12},&
 \gamma_j&=-\frac{3\cos2\alpha_j}{\sin\alpha_j\sin3\alpha_j},\\
 \Delta X'_1&=-\frac{\log A_{12}}{k_1}+\gamma_1M_2,&
 \Delta X'_2&=\frac{\log A_{12}}{k_2}-\gamma_2M_1.
 \label{eq:shifts}
\end{align}
Direct differentiation of the root-cross-ratio expression verifies
\(\partial_{P_j}\delta_0=-\Delta X'_j\), without defining one side
from the other. It also verifies antisymmetry of the derived responses
under simultaneous exchange of angles and species.

\paragraph{Conditional interface \Cc/\Cd\ (\ct, \grid).}
First-order inversion of
\(P_j^{\rm phys}=P_j(\alpha_j)+\hbar p_j(\alpha_j)+O(\hbar^2)\)
gives the algebraic conversion
\begin{equation}
 \delta_{1,P}=\delta_{1,\alpha}+R[p],\qquad
 R[p]=\sum_j\Delta X'_jp_j.
 \label{eq:conversion}
\end{equation}
The physical identification of both \(p_j\) and
\(\delta_{1,\alpha}\) is still \Cd. The scripts retain two audited
trace representatives to expose this issue:
\begin{align}
 p_{\rm m}(\alpha,\epsilon)
 &=-\frac{\epsilon}{2\pi}\sum_r c_r\sin\varphi_r
                 \log|\tan(\varphi_r/2)|,\\
 e_{\rm m}(\alpha)
 &=\frac{-6\sin\alpha-2\sum_r c_r\sin\varphi_r\cos\varphi_r
                    \log|\tan(\varphi_r/2)|}{4\pi},\\
 L_1&=vp_{\rm m}-e_{\rm m},\qquad p_{\rm can}=\frac{\dd L_1}{\dd v},\\
 (\varphi_1,\varphi_2,\varphi_3)
 &=(\alpha+2\pi/3,\alpha-2\pi/3,\alpha),\quad(c_1,c_2,c_3)=(1,1,-2).
 \label{eq:charges}
\end{align}
\Cc The mode momentum is checked by independent half-line integration.
The canonical representative is the derivative of the same conditional
trace, not a new physical momentum determination. Switching from
\(p_{\rm m}\) to \(p_{\rm can}\) requires changing the label-frame
phase by \(-R[p_{\rm can}-p_{\rm m}]\); otherwise it changes the answer.
Agreement of the fixed-momentum one-body energy does not justify that
inconsistent switch in a two-body phase.

\section{Elastic and other bulk counterterm responses}
\paragraph{Exact rescaling \Cc\ (\ct).}
Consider the family, with positive \(a,b,d\) near one,
\begin{equation}
 H_{a,b,d,g}=\int\left(\frac a2\pi^2+\frac b2u^2-gu^3+
                               \frac d2u_x^2\right)\dd x.
\end{equation}
The transformation
\begin{equation}
 x=L_xX,\quad t=L_tT,\quad \phi=F\Phi,\qquad
 L_x=\sqrt{d/b},\quad L_t=\frac{\sqrt{d/a}}b,\quad F=\frac{\sqrt{bd}}g
 \label{eq:scaling}
\end{equation}
maps its action to \(s\) times the fixed action and its momentum to
\(p_*\) times the fixed momentum, where
\begin{equation}
 s=\frac{db^{3/2}}{\sqrt a\,g^2},\qquad
 p_*=\frac{b^2\sqrt d}{\sqrt a\,g^2},\qquad
 \delta_{a,b,d,g}(P_1,P_2)
 =s\,\delta_0(P_1/p_*,P_2/p_*).
 \label{eq:scaledphase}
\end{equation}
The last expression uses the continuously rescaled classical Darboux
coordinates and fixed laboratory momentum units. Every transformed
Lagrangian coefficient and the independently transformed canonical
momentum are checked symbolically. The crucial momentum factor for an
elastic change is \(b^2\), whereas the action factor is \(b^{3/2}\).

\Cc Set
\[
 a=1+2\kappa\hbar,\quad b=1+2\lambda\hbar,\quad
 g=1-\eta\hbar,\quad d=1+2\zeta\hbar.
\]
Thus \(\Delta H/\hbar=\int(\lambda u^2+\eta u^3+
\zeta u_x^2+\kappa\pi^2)\dd x\). With
\(D=\sum_jP_j\Delta X'_j\), expansion of the exact rescaling gives
\begin{align}
 K_\lambda&=3\delta_0+4D,& K_\eta&=2\delta_0+2D,\\
 K_\zeta&=2\delta_0+D,& K_\kappa&=-\delta_0-D.
 \label{eq:responses}
\end{align}
All four are checked by numerically solving the exact deformed momentum
equations before taking a central difference. At 75 digits and step
\(10^{-18}\), the largest discrepancy is \(1.87\times10^{-35}\),
below the unchanged \(10^{-32}\) assertion. This is a finite-difference
validation error, not an uncertainty in an unknown quantum determinant.

\Cc In particular,
\begin{equation}
 \boxed{K_\lambda=\frac23(v_2-v_1)\log A_{12}
       +3T_{12}+4(P_1\gamma_1M_2-P_2\gamma_2M_1).}
 \label{eq:lambda}
\end{equation}
The response changes sign under complete label exchange; hence its
symmetric part vanishes. At \((\alpha_1,\alpha_2)=(0.3,0.6)\),
\[
 K_\lambda^{++}=0.05008140486841890958\ldots,\qquad
 K_\lambda^{+-}=-0.22517393868334175848\ldots.
\]
Consequently there is no general \(\lambda\)-independence of the
antisymmetric canonical phase. The statement is for the specified
classical-frame response, with other bulk and boundary couplings held
fixed. It is not a claim that an uncomputed quantum boundary prescription
is independent of all possible scheme changes.

\paragraph{Independent connected insertion \Cc\ (\insertcheck).}
The bilinear equation implies, for \(f=\log\tau\),
\(f_{tt}-f_{xx}+f_{xxxx}+6f_{xx}^2=0\).
Let \(F(t)\) be the spatial integral of
\(\log\tau_{12}-\log\tau_1-\log\tau_2\), with the time-independent
right plateau \(\log A_{12}\) subtracted. Then
\begin{equation}
 \int\dd t\dd x\,[u_{12}^2-u_1^2-u_2^2]
 =-\frac23[F'(+\infty)-F'(-\infty)]
 =\frac23(v_1-v_2)\log A_{12}.
 \label{eq:insertion}
\end{equation}
Indeed \(F'(-\infty)=-v_2\log A_{12}\) and
\(F'(+\infty)=-v_1\log A_{12}\). Independent spatial quadratures
at \(t=\pm3000\) test these slopes in both channels. The minus sign
from \(\Delta L=-\lambda\hbar u^2\) reproduces the logarithmic
part of \eqref{eq:lambda}. The remaining finite part is the difference
between the connected bulk insertion and the response of the canonical
scattering generator, including its changing endpoint/Darboux terms.
Dropping those terms would not implement the stipulated fixed canonical
momenta and frame. Their quantum state interpretation is not resolved here.

\paragraph{What ``other allowed counterterms'' means.}
\Cc\ (\ct, \three) A finite tadpole \(\mu\hbar\int u\dd x\)
has zero connected scattering insertion: the exact tau asymptotes give
\(M_N=\sum_j2k_j\) at every time. Vacuum constants also cancel after
vacuum normalization. For a uniform background \(U\), the large-momentum
expansion of the one-loop energy integrand is checked directly from
\begin{equation}
 \tfrac12\sqrt{q^2+q^4-6Uq^2}
 =\tfrac12|q|\sqrt{1+q^2}
 -\frac{3U|q|}{2\sqrt{1+q^2}}
 -\frac{9U^2|q|}{4(1+q^2)^{3/2}}+O(U^3).
\end{equation}
Only its term linear in \(U\) is UV divergent in one spatial dimension.
For a one-loop background ring with \(n\) insertions the superficial
anisotropic degree is \(3-2n\): each insertion contributes at most
two internal spatial derivatives, each propagator degree \(-4\),
and the frequency/momentum measure degree \(3\). Thus no further
bulk divergent counterterm is forced at this order in the specified
cubic theory. These are perturbative UV statements, not infrared or
boundary renormalisation results.

\Cd Symmetries alone do not restrict finite local modifications to
\eqref{eq:responses}. Terms such as \(\hbar c_nu^n\) and higher
derivative terms are possible changes of the theory; boundary operators
are a separate issue. Their finite coefficients are not fixed by the
given data. We parameterise their possible contribution as
\(\sum_A c_AK_A\), with \(K_A\) \emph{uncomputed} outside the
four-coupling family above, rather than assert an exhaustive finite
classification. No new higher interaction is silently set to a preferred
value. The response \(K_\lambda\) is ready to accept the parallel
renormalisation package's coefficient without choosing that coefficient.

\section{Continuum integrals and the symmetric/antisymmetric split}
\paragraph{Definitions and independent checks \Cc\ (\kernels, \grid).}
Use the Jost phase continuous on each open half-line and vanishing at
infinity. With the roots and weights of \eqref{eq:charges},
\begin{align}
 \theta(q;\alpha,\epsilon)=-\sum_r c_r\bigg[&
 \operatorname{sgn}(q)\arctan\frac{\sin\varphi_r}{|q|}\\[-1ex]
 &+\epsilon\arctan\frac{\sin\varphi_r\cos\varphi_r}
 {|q|\sqrt{1+q^2}+q^2+\sin^2\varphi_r}\bigg],\nonumber\\
 \theta'(q)&=\sum_r\frac{c_r\sin\varphi_r}{q^2+\sin^2\varphi_r}
 \left(1+\frac{\epsilon\operatorname{sgn}(q)\cos\varphi_r}{\sqrt{1+q^2}}\right).
\end{align}
The transmission obtained directly from the KP-root cross ratio agrees
with \(e^{i\theta}\). Differentiation of the arctangent formula agrees
with the independently written derivative. Define ordinary half-line
integrals, with no delta function inserted at zero,
\begin{equation}
 I_{12}=\frac1{4\pi}\int_{\mathbb R\setminus\{0\}}
       \theta_1\theta'_2\dd q,\qquad
 J_{12}=\frac{I_{12}-I_{21}}2,\qquad
 I_{12}^{\rm sym}=-\frac\pi4(\epsilon_1+\epsilon_2).
 \label{eq:integrals}
\end{equation}
The last equality follows from the endpoints
\(\theta(0^\pm)=\pi(\pm1+\epsilon)\). Independent quadrature for
all four orientations gives maximum closed-form discrepancy
\(5.94\times10^{-66}\) at 65 digits.

\Cc For reproducibility the closed form uses
\begin{align}
 L(z)&=\int_0^z\frac{2\log|t|}{1-t^2}\dd t,\qquad
 L(1)=-\pi^2/4,\quad L(-z)=-L(z),\\
 I^{++}(\alpha_1,\alpha_2)
 &=\frac1{2\pi}\sum_{rs}c_rc_sL\left(
       \frac{\tan(\varphi_{1r}/2)}{\tan(\varphi_{2s}/2)}\right),\\
 I^{+-}(\alpha_1,\alpha_2)
 &=\frac1{2\pi}\sum_{rs}c_rc_sL\left(
       \tan(\varphi_{1r}/2)\tan(\varphi_{2s}/2)\right),\\
 I^{\epsilon_1\epsilon_2}
 &=\tfrac12(\epsilon_1+\epsilon_2)I^{++}
   +\tfrac12(\epsilon_1-\epsilon_2)I^{+-}.
\end{align}
These are real half-line formulas, not a prescription for holomorphic
crossing beyond their stated domain. In particular \(I\) and \(J\)
are alternative representatives sharing a local gradient, not additive
pieces of the phase. The raw symmetric term is an endpoint artifact;
it is not a physical symmetric scattering phase established by this work.

\paragraph{Exactly what is independent of \(\lambda\).}
\Cc At one-loop order the displayed undeformed-background integrals
\(I,J\), their endpoint symmetric term, and the reference-trace
conversion \(R[p_{\rm m}]\) are \(\lambda\)-independent. An
order-\(\hbar\) deformation of their arguments affects the scattering
phase only at the next loop order. All explicitly computed local
\(\lambda\) dependence at this order is \(\lambda K_\lambda\)
in the rescaled classical frame. The antisymmetric \(J\) alone is
independent of \(\lambda\); the antisymmetric canonical response is
not. Nothing here proves scheme independence of the unknown boundary
contribution under more general changes of the theory.

\Cd A useful but conditional parametrisation is
\begin{equation}
 \delta_{1,P}=J_{12}+R[p_{\rm m}]+\lambda K_\lambda
 +\eta K_\eta+\zeta K_\zeta+\kappa K_\kappa+B_{12}
 +\sum_Ac_AK_A.
 \label{eq:unknown}
\end{equation}
Here \(B_{12}\) includes the missing boundary/collective determinant,
state map, and absolute normalization relative to the displayed trace
representative. It is not known to be a constant. If scalar exchange is
identified with inverse physical scattering, a locally continuous unitary
representative has symmetric part zero, and \(B_{21}=-B_{12}\) modulo
the phase branch. This algebraic parametrisation does not prove that the
quantum theory is elastic or factorised.

\begin{center}\small
\begin{tabular}{@{}ccrrrr@{}}\toprule
Channel & $(\alpha_1,\alpha_2)$ & $J$ & $R[p_{\rm m}]$ & $R[p_{\rm can}]$ & $K_\lambda$\\\midrule
$++$ & $(.2,.4)$ & .412598425 & 2.479485715 & 3.699829027 & .064969423\\
$++$ & $(.3,.6)$ & .392706352 & 1.968010544 & 3.666344128 & .050081405\\
$++$ & $(.4,.6)$ & .227909672 & .811937649 & 2.079893527 & .057940143\\
$+-$ & $(.2,.4)$ & .217194901 & 4.825323137 & 6.747182050 & .061623021\\
$+-$ & $(.3,.6)$ & .408542241 & 4.915210065 & 8.209382585 & -.225173939\\
$+-$ & $(.4,.6)$ & .478642156 & 5.241126358 & 9.896087545 & -1.073366855\\\bottomrule
\end{tabular}
\end{center}
\Cc\ (\grid) These are ingredients, not two competing physical
answers obtained by selecting a preferred momentum convention.
For example \(J+R[p_{\rm m}]=2.360716896635\ldots\) at
\((.3,.6,++)\); its difference from the full phase has no numerical
error bound because \(B_{12}\) is uncomputed.

\section{Three-soliton test: what factorisation does and does not fix}
\paragraph{Exact mode construction \Cc\ (\three).}
For \(E_j=\exp[k_j(x-v_jt-x_{0j})]\), define
\begin{equation}
 \tau_N(E)=\sum_{\mu\in\{0,1\}^N}
        \prod_j E_j^{\mu_j}\prod_{i<j}A_{ij}^{\mu_i\mu_j}.
\end{equation}
Adding a fourth infinitesimal wave \(K=iq\),
\(W=i|q|\sqrt{1+q^2}\) to \(N=3\) gives
\begin{equation}
 \delta\phi=2\partial_x\left[
     e^{iqx-i\omega t}\frac{\tau_3(T_1E_1,T_2E_2,T_3E_3)}
                                  {\tau_3(E_1,E_2,E_3)}\right],
 \qquad T_j=A(j,\mathrm{ph}),\quad \omega=|q|\sqrt{1+q^2}.
 \label{eq:threephonon}
\end{equation}
The coefficient test expands the full four-wave bilinear equation using
\(\mathcal P(K,W)=W^2-K^2+K^4\) and
\(A_{ij}=-\mathcal P(K_i-K_j,W_i-W_j)/
\mathcal P(K_i+K_j,W_i+W_j)\). It checks all 81 exponential
coefficients exactly for a rational three-soliton configuration and one
imaginary phonon. Independent 75-digit checks cover three right-movers
and a mixed-direction triple, with both signs of phonon momentum.
This is explicit coefficient evidence, not a spectral completeness proof.

\Cc Exact polynomial clustering gives
\begin{align}
 \tau_3|_{E_3=0}&=\tau_2(E_1,E_2),\\
 \lim_{E_3\to\infty}\frac{\tau_3}{E_3}
 &=\tau_2(A_{13}E_1,A_{23}E_2),\\
 \lim_{E_3\to\infty}
 \frac{\tau_3(T_jE_j)}{\tau_3(E_j)}
 &=T_3\frac{\tau_2(T_1A_{13}E_1,T_2A_{23}E_2)}
                 {\tau_2(A_{13}E_1,A_{23}E_2)}.
\end{align}
Thus the continuum transmission is \(\prod_jT_j\). In a sequence of
separated binary collisions, the \emph{formal continuum representative}
is the sum of the oriented \(J_{ij}\). This is what the script records;
the full three-body determinant entry is \texttt{null}.

\paragraph{A precise insufficiency result \Cc\ (\three).}
Suppose a scalar three-body amplitude has the stipulated cluster form
\(S_3=S_{12}S_{13}S_{23}\). For any real antisymmetric pair function
\(F_{ij}\), the replacements
\begin{equation}
 S_{ij}\longmapsto e^{iF_{ij}}S_{ij},\qquad
 S_3\longmapsto e^{i(F_{12}+F_{13}+F_{23})}S_3
 \label{eq:pairfreedom}
\end{equation}
preserve that factorisation and scalar collision-order consistency.
An explicit example is
\begin{equation}
 F_{12}=c(P_1M_2-P_2M_1),\qquad c\in\mathbb R.
 \label{eq:F}
\end{equation}
It is nonzero, antisymmetric, vanishes on a zero-charge leg, and has
rank at most two as a kernel of one-body charges. Its \(3\times3\)
minors vanish identically. It also changes sign under the formal map
\((P,M)\mapsto(-P,-M)\) on either leg. This last algebraic property
is not an asserted physical crossing axiom for the good theory.

The example does not evade the Toda audit: that audit additionally has
a microscopic Fredholm/Krein equation, and its excluded example has
infinite rank. Neither the finite-rank condition nor scalar clustering
by itself eliminates \eqref{eq:F}. A fully determined Boussinesq
determinant gradient would eliminate additions that change that gradient;
the point is that such a gradient has not been established on the
physical constrained domain.

\Cd Thus computing the exact three-soliton \emph{phonon modes} and
imposing clustering does not compute its one-loop \emph{phase}, and
does not fix \(B_{12}\) or the residual \(\vartheta\).
This result concerns sufficiency of the proposed test, not a proof of
arbitrary physical freedom. Spatial removal of a spectator at finite
observation time must also be distinguished from an infinite-time
scattering experiment in which that spectator eventually collides.

\paragraph{The boundary obstruction persists \Cc/\Cd\ (\three).}
For the centered potential \(\Phi=k\tanh(kx/2)\),
\begin{equation}
 \Phi_\alpha=\cos\alpha\,[\tanh(kx/2)+(kx/2)\operatorname{sech}^2(kx/2)],
 \qquad \|\Phi_\alpha\|_{[-R,R]}^2=2R\cos^2\alpha+O(1).
\end{equation}
The tail slope is rechecked by independent quadrature. The controlling
projector audit gives localized data whose collective projection acquires
a nonzero endpoint tail. Those collective and radiation directions cannot
be integrated independently with both held in a fixed-endpoint domain.
Even at finite dimension,
\(\det\left(\begin{smallmatrix}A&B\\B^T&C\end{smallmatrix}\right)
=\det C\det(A-BC^{-1}B^T)\); knowledge of the radiation block does not
compute the Schur factor. The scripts check this identity in a nontrivial
scalar-block example. \Cd A common domain, propagated constraints,
completeness, the Schur factor, and temporal state normalization remain
necessary. No full determinant or relative Maslov index is claimed here.

\section{Approach to resonance and cached data}
\paragraph{Regular head-on Y-boundary \Cc\ (\grid).}
For \(\epsilon_1=+1,\epsilon_2=-1\), fix
\(0<\alpha_1<\pi/3\) and put
\(\alpha_2=\pi/3-\alpha_1-h\), \(h\downarrow0\).
The root denominator vanishes and \(A_{12}\to+\infty\).
The integrals have a finite limit, checked by direct endpoint quadrature.
For either specified momentum function \(p\),
\begin{align}
 J+R[p]&=C_p\log A_{12}+B_p+o(1),\\
 C_p&=\left[-p_1/k_1+p_2/k_2\right]_{h=0},\qquad
 B_p=\left[J+\gamma_1M_2p_1-\gamma_2M_1p_2\right]_{h=0},\\
 K_\lambda&=C_\lambda\log A_{12}+B_\lambda+o(1),\qquad
 C_\lambda=\left[\tfrac23(v_2-v_1)\right]_{h=0},\\
 B_\lambda&=\left[3T_{12}
       +4(P_1\gamma_1M_2-P_2\gamma_2M_1)\right]_{h=0}.
 \label{eq:resonance}
\end{align}
At \(\alpha_1=0.4\), the checked numbers are
\begin{align*}
 J_*&=0.516497248493590315982850593057\ldots,\\
 C_{p_{\rm m}}&=1.27122424749067227378534580175\ldots,\\
 C_\lambda&=-1.14589177876216227606321733046\ldots.
\end{align*}
Therefore the logarithmic coefficient of the conditional mode
representative is \(C_{p_{\rm m}}+\lambda C_\lambda\) when the
other computed bulk changes are absent. \Cd The unknown boundary
phase can change the full asymptotics; this is not a prediction of a
physical resonant amplitude or fusion rule. The expansion is taken at
fixed regular angles before the resonance limit and is not asserted
uniform when \(\hbar|\log h|\) becomes order one.

\Cc\ (\grid) The genuine overtaking Y-line is
\(\alpha_1+\alpha_2=2\pi/3\). It has no interior point with both
legs on the light branch. We save classical responses approaching this
line with one heavy leg, with quantum phases explicitly absent.
The light overtaking \(A\to0\) limit on the diagonal is instead an
equal-velocity threshold; the sampled \(K_\lambda\) tends to zero,
not to a nonzero logarithmic divergence. The historical
\(|\alpha_1-\alpha_2|=\pi/3\) overtaking resonance is false:
\(A(0.2,0.2+\pi/3)=0.355465479849\ldots\), finite and nonzero.
The grid does not analytically continue a real one-loop phase through
an unstable heavy branch or a singular \(A<0\) tau function.

\paragraph{Files and error budget \Cc\ (all five scripts).}
\file{good_twobody_grid.json} covers the Cartesian angles
\(0.1,0.2,\ldots,1.0\), excluding equal velocities and nonpositive
head-on \(A\), in all four orientations: 270 rows.
Each row is recalculated at 55 and 75 digits; the largest difference
in its tested numerical fields is \(8.31\times10^{-55}\).
The independent quadrature tolerance is \(10^{-55}\), the fixed-momentum
finite-difference tolerance \(10^{-32}\), and tabulated decimals are
rounded. These are deterministic numerical consistency estimates, not
rigorous interval enclosures. They do not cover the unknown physical
determinant, whose uncertainty is unbounded by this calculation.

\file{good_twobody_resonance.json} contains three head-on approaches
at \(\alpha_1=0.2,0.4,0.5\) with
\(h=10^{-2},10^{-4},\ldots,10^{-12}\), limiting coefficients and
subtracted remainders, plus the overtaking threshold and heavy Y-line
diagnostics. At the last head-on point the maximum discrepancy from the
subtracted limits is below \(10^{-8}\); individual finite-distance
errors are stored rather than identified with quadrature errors.
The other four JSON files hold kernel, scaling, mode, and local-insertion
verification results. Strings retain 55 significant digits; \texttt{null}
means uncomputed, never zero.

\clearpage
\section{Claim/check summary and failures recorded}
\small
\begin{longtable}{@{}>{\raggedright\arraybackslash}p{8cm}>{\raggedright\arraybackslash}p{6.4cm}@{}}\toprule
Claim and status & New checker\\\midrule\endhead
\Cc Continuum $I,J$, all orientations; endpoint symmetric term & \kernels; \grid\\
\Cc Fixed-momentum response to $u^2,u^3,u_x^2,\pi^2$; antisymmetry & \ct\\
\Cc Connected $u^2$ insertion and canonical endpoint difference & \insertcheck\\
\Cc Four-wave coefficient tests; exact three-soliton clustering & \three\\
\Cc Scalar clustering is insufficient to fix the pair phase & \three, rank-two example\\
\Cc Regular grid and qualified resonance asymptotics & \grid\\
\Cd Complete projected determinant, absolute physical phase, quantum state map & Not computed; explicit unknown entries\\
\Cd Other finite local interactions and quantum resonant channels & Not determined by the supplied theory and checks\\\bottomrule
\end{longtable}
\normalsize

\paragraph{Failures recorded.}
\begin{enumerate}
\item The first connected-insertion quadrature subtracted three nearly
constant temporal log-tau derivatives numerically on an infinite interval.
Its unchanged \(10^{-35}\) slope test failed at the overtaking point,
\(t=-3000\), with error \(8.2718016287\times10^{-4}\).
The plateau cancellation was performed algebraically before integration,
and equality of the original and stabilized integrands was proved with
SymPy. The same test and tolerance were retained. This failure is not
hidden by the final success markers.
\item An initial hand scaling assigned momentum the action factor
\(b^{3/2}\). Direct transformation of
\(P=-\int\pi\phi_x\dd x\) corrected it to \(b^2\) before
the numerical checks. The exact transformed momentum and independent
inverse-momentum tests are included to prevent recurrence.
\item The intended inference ``factorised three-soliton phonons plus
clustering fixes the quantum antisymmetric phase'' fails as a logical
completion strategy: \eqref{eq:pairfreedom}--\eqref{eq:F} preserve the
tested constraints. This is a negative result, not a failed numerical
test repaired by deleting a constraint. The full three-body phase was
not computed.
\item The inherited identification of a continuum integral with the
absolute determinant, the old false overtaking resonance line, and the
uncorrected classical energy are not adopted. Their controlling audits
and errata retain precedence. No check was weakened and no legacy
failure was claimed to have been repaired by this package.
\item The first PDF build had two overwide equation lines and one
overwide command line. They were split; an orphan section heading and
a split reading-scope paragraph were moved to fresh pages before the
final two-pass build and rendered inspection. Scientific formulas were unchanged.
\end{enumerate}

\paragraph{Verification commands.}
Run each of the five named scripts with the environment variable
\texttt{PYTHONDONTWRITEBYTECODE=1} and command
\texttt{python3 phase2/scripts/}\emph{script}.
Each successful run ends with \texttt{ALL CHECKS PASSED} only after its
assertions. Total runtime is comfortably below 30 minutes on this machine;
the 270-row grid took about 107 seconds. Build this note by running
\texttt{pdflatex -interaction=nonstopmode -halt-on-error good\_twobody.tex}
twice in \file{phase2/}. No project-wide script loop is needed.

\clearpage
\section*{Reference verification and reading scope}
The repository inputs listed in Section 1 were read, with particular
attention to the audit verdicts. For the external reference below, the
local PDF title page and relevant full-text passages on printed pages
52--53, 62 and 65--66 were read; the entire paper was \emph{not} read
for this package. No new arXiv identifier or DOI is asserted.
\Ca Faddeev--Korepin's discussion explicitly concerns a Klein--Gordon
fluctuation determinant and leaves a factor \(Z\) to supplementary
requirements in its two-body formula. \Cd Applying its full determinant
argument to the constrained fourth-order Boussinesq problem remains
unproved here; the reference is not cited as that missing theorem.

\begin{thebibliography}{1}
\bibitem{FK} L.~D.~Faddeev and V.~E.~Korepin,
\emph{Quantum theory of solitons}, Physics Reports \textbf{42}, no.~1
(1978), 1--87. [Local PDF, title page and specified full-text passages read.]
\end{thebibliography}
\end{document}
